%% kapa: exam=1 round=37 year=2026 period=1 form=A subject=civil_law subject_ko=민법 seq=20 number=20
%% source: Q-Net 감정평가사 기출문제 (cst00309) artlSeq=5258924 file=raw/1008/5258924/2026년 제37회 감정평가사 1차 1교시.pdf
%% conversion: kordoc 4.24.51 -> tools/gen1.py (auto); review=proofread-sonnet; flags=none
\begin{question}{20}
소멸시효의 중단에 관한 설명으로 옳은 것을 모두 고른 것은? (다툼이 있으면 판례에 따름)
\begin{framedtext}
ㄱ. 여러 차례의 최고가 있은 후, 재판상 청구가 있더라도 그 시효중단의 효력은 항상 최초의 최고를 한 때에 발생한다. \par
ㄴ. 소멸시효의 중단사유로서의 승인은 소멸시효의 진행이 개시된 이후에만 가능하고, 현존하지 아니하는 장래의 채권을 미리 승인하는 것은 채무자가 그 권리의 존재를 인식하고서 한 것이라고 볼 수 없어 허용되지 않는다. \par
ㄷ. 물상보증인이 그 피담보채무의 소멸을 이유로 제기한 저당권설정등기 말소등기절차이행청구소송에서 채권자 겸 저당권자가 청구기각의 판결을 구하고 피담보채권의 존재를 주장하였다고 하더라도 피담보채권의 소멸시효에 관하여 규정한 민법 제168조 제1호 소정의 ‘청구’에 해당하지 아니한다. \par
\end{framedtext}
\begin{choices}
  \item ㄱ
  \item ㄴ
  \item ㄱ, ㄷ
  \item ㄴ, ㄷ
  \item ㄱ, ㄴ, ㄷ
\end{choices}
\end{question}
\answer{4}
